Reports documenting the failure modes of interpretability methods are frequent and highlight a recurring theme: a general statistical fragility, most strikingly illustrated by dead salmon artifacts. We provide here a non-exhaustive overview of such issues.

\medskip
\noindent\textbf{Feature Attribution.} 
Gradient-based attribution methods \citep{simonyan2014deepinsideconvolutionalnetworks,sundararajan2017axiomatic} aim to highlight input features most relevant to model predictions. However, \citet{adebayo2018sanity} demonstrated that saliency maps can remain visually plausible even after model weights are randomized. Further, \citet{NEURIPS2019_bb836c01} showed that gradient-based explanations can be manipulated by adversarial perturbations, leaving predictions unchanged, while \citet{Ghorbani_Abid_Zou_2019} revealed that explanations are unstable under minor data transformations. From a theoretical standpoint, \citet{bilodeau2024impossibility} established impossibility results showing that no attribution method can simultaneously satisfy intuitive desiderata across broad model classes. 

\medskip
\noindent\textbf{Probing.} 
Probing methods train a classifier to predict a target label from internal activations. Early studies already showed that both linear and structural probes could recover information with surprisingly high accuracy from randomized contextualized embeddings \citep{conneau-etal-2018-cram,hewitt-manning-2019-structural}, and syntactic probes do not generalize~\citep{hall-maudslay-cotterell-2021-syntactic}. Later, \citet{ravichander-etal-2021-probing} demonstrated that probes can extract features merely encoded (e.g., inherited from embeddings) even if unused during inference; probing asks whether a concept is encoded in an activation, not whether it is computationally relevant. Capacity-controlled and information-theoretic probes~\citep{voita-titov-2020-information,zhu-rudzicz-2020-information,pimentel-etal-2020-information,belinkov-2022-probing} or amnesic probing~\citep{elazar-etal-2021-amnesic} attempt to mitigate such false discoveries. 

\medskip
\noindent\textbf{Sparse Autoencoders.} 
Unsupervised concept-discovery pipelines such as sparse autoencoders (SAEs) \citep{cunningham2023sparseautoencodershighlyinterpretable,yun-etal-2021-transformer,bricken2023monosemanticity,templeton2024scaling} display analogous pathologies. 
\citet{heap2025sparseautoencodersinterpretrandomly} showed that SAEs can recover apparently interpretable components even in randomly initialized transformers. Additional studies show that SAEs often fail to generalize across settings or tasks 
\citep{heindrich2025sparseautoencodersgeneralizecase,kantamneni2025sparseautoencodersusefulcase}. Also, \citet{li2025interpretabilityillusionssparseautoencoders} show SAE are sensitive to adversarial input perturbations.

\medskip
\noindent\textbf{Concept-Based Explanations.}
Concept-based methods \citep{kim2018interpretability,bau2017network} aim to identify human-interpretable concepts that align with model representations (e.g., concept activation vectors~\citep{kim2018interpretability} or network dissection~\citep{bau2017network}). These methods also face documented limitations~\citep{sinha2025comprehensivesurveyriskslimitations,aysel2025conceptbasedexplainableartificialintelligence}. Already, \citet{bolukbasi2021interpretabilityillusionbert} showed \textit{interpretability illusion} arising where activations of individual neurons in BERT may spuriously appear to encode a concept. Then, \citet{nicolson2025explainingexplainabilityrecommendationseffective} showed that concept activation scores can produce highly inconsistent explanations, and \citet{ramaswamy2023overlooked} documented poor generalization and high sensitivity to the dataset used to infer concepts. Finally, \citet{piratla2024estimationconceptexplanationsuncertainty} further demonstrated high variance and recommended incorporating uncertainty estimation.


\medskip
\noindent\textbf{Causal Approaches.}
To address issues with prediction-based explanations, a shift toward causality-based interpretability has emerged through the use of causal mediation analysis \citep{pearl2012causal,elazar-etal-2021-amnesic,NEURIPS2020_92650b2e,meng2022locating,finlayson-etal-2021-causal,syed2023attribution,monea2024glitchmatrixlocatingdetecting,mueller2024questrightmediatorhistory}. These methods intervene on intermediate representations to quantify causal effects of components on model outputs. 
Yet recent work documents substantial fragilities and trade-offs~\citep{canby2025reliablecausalprobinginterventions}: \citet{zhang2024bestpracticesactivationpatching} showed that such approaches are sensitive to experimental design. Then, \citet{mcgrath2023hydraeffectemergentselfrepair} discovered a ``hydra effect,'' where ablating components identified as causally important fail to change behavior due to redundant causal pathways. This phenomenon, known as \textbf{overdetermination}, occurs when multiple redundant, independently sufficient causal pathways exist~\citep{schaffer2003overdetermining,sider2003s,Dyrkolbotn_2017}. Rather than isolating simple mechanisms, interventions tend to reveal overdetermined causal structures.
 
\medskip
\noindent\textbf{Mechanistic Interpretability.}
Causal approaches culminate in \textit{mechanistic interpretability} (MI), which aims to reverse-engineer networks into human-interpretable algorithms \citep{olah2020zoom}. One family of approaches (\textit{where-then-what}) first identifies circuits carrying information from inputs to outputs and then interprets their components \citep{dunefsky2024transcodersinterpretablellmfeature,davies2024cognitiverevolutioninterpretabilityexplaining,NEURIPS2023_34e1dbe9}. The second (\textit{what-then-where}) instead starts from high-level candidate algorithms and searches for causally aligned neural subspaces, using \textit{causal abstraction} metrics \citep{Geiger-etal:2022:SAIL,pmlr-v162-geiger22a,Beckers_Halpern_2019}. Despite promising demonstrations, both categories have the typical issues~\citep{sharkey2025openproblemsmechanisticinterpretability}. Subspace patching can produce \textit{interpretability illusions} by activating alternate pathways~\citep{makelov2023subspacelookingforinterpretability}, also a problem of overdetermination. Circuit explanations often fail to generalize \citep{wang2022interpretabilitywildcircuitindirect,li2025interpretationpredictbehaviorunseen} and are sensitive to minor experimental choices \citep{meloux2025mechanisticinterpretabilitystatisticalestimation}. Exhaustive studies on toy models reveal multiple incompatible explanations for both strategies, even for random networks \citep{melouxeverything}. Finally, \citet{sutter2025nonlinearrepresentationdilemmacausal} proved that, in general, existing causal abstraction methods can produce explanations for random networks. 


\medskip
\noindent\textbf{Natural Language Explanations.}
Generating natural language rationales has become a popular interpretability approach \citep{marasovic-etal-2022-shot,wiegreffe-etal-2022-reframing}. However, \citet{ajwani2024llmgeneratedblackboxexplanationsadversarially} showed that LLM-generated explanations can be systematically unfaithful, confidently providing plausible-sounding justifications for predictions made for entirely different reasons. Moreover, chain-of-thought (self-)explanations are typically unfaithful to the model’s computation~\citep{lanham2023measuringfaithfulnesschainofthoughtreasoning,arcuschin2025chainofthought,turpin2023languagemodelsdontsay}. There exist infinite plausible stories that can rationalize any behavior post hoc. Thus, natural language explanations are particularly susceptible to confabulation and, thus, to false positives. 